%%%%%%%%%%%%%%%%%%%%%%%%%%%%%%%%%%%%%%%%
%
% $ Autor: Gruppe 07 $
% $ Projekt: Magic Faucet Fountain $
% $ Pfad: LaTeXTemplate/Montageanleitung/Chapters/de/Funktionstest.tex $
% $ Kurzbeschreibung: Dieses Dokument erstellt die Inbetriebnahme der Montageanleitung $
%
% !TeX spellcheck = de_DE/GB
% !TeX program = pdflatex
% !BIB program = biber/bibtex
% !TeX encoding = utf8
%
%%%%%%%%%%%%%%%%%%%%%%%%%%%%%%%%%%%%%%%%

\chapter{Inbetriebnahme und Funktionstest}
Nach dem vollständigen Aufbau der Baugruppe \textbf{Magic Faucet Fountain} kann die Inbetriebnahme erfolgen. Diese umfasst die Aktivierung der Elektronik, die Verbindung mit der Smartphone-App sowie die Überprüfung des Wasserflusses.

\subsubsection{Vorbereitung}
\begin{itemize}
	\item Überprüfen Sie, ob alle Komponenten korrekt montiert und fest miteinander verbunden sind.
	\item Achten Sie darauf, dass der Wasserbehälter ausreichend befüllt ist und keine Undichtigkeiten vorhanden sind.
	\item Vergewissern Sie sich, dass die Pumpe vollständig im Wasser eingetaucht ist.
\end{itemize}

\subsubsection{Inbetriebnahme}
\begin{enumerate}
	\item \textbf{Anschluss an das Stromnetz:} Stecken Sie das Netzkabel des Elektronikgehäuses in eine Steckdose.
	\item \textbf{LED-Anzeige prüfen:}
	\begin{itemize}
		\item Eine \textbf{grüne LED} leuchtet dauerhaft – dies signalisiert, dass die Baugruppe mit Strom versorgt wird.
		\item Eine \textbf{blaue LED blinkt} – dies zeigt an, dass das Bluetooth-Modul bereit zur Verbindung ist.
	\end{itemize}
	\item \textbf{Bluetooth-Verbindung herstellen:}
	\begin{itemize}
		\item Öffnen Sie die nRF Connect-App.
		\item Suchen Sie das Gerät in der Liste verfügbarer Bluetooth-Verbindungen.
		\item Nach erfolgreicher Verbindung leuchtet die \textbf{blaue LED dauerhaft}.
	\end{itemize}
	\item \textbf{Wasserfluss aktivieren:}
	\begin{itemize}
		\item Betätigen Sie den \textbf{Kippschalter} am E-Kasten.
		\item In der App kann der Wasserfluss nun zusätzlich manuell ein- oder ausgeschaltet werden.
		\item Die App zeigt zudem den aktuellen Wasserstand im System an.
	\end{itemize}
\end{enumerate}

\subsubsection{Funktionstest}
\begin{itemize}
	\item Nach dem Einschalten der Pumpe sollte Wasser gleichmäßig durch das Rohr zum schwebenden Hahn befördert werden.
	\item Der austretende Wasserstrahl sollte das Rohr vollständig verdecken, sodass der gewünschte Effekt entsteht.
	\item Über die App lässt sich der Wasserfluss jederzeit steuern.
	\item Kontrollieren Sie die LEDs auf dem E-Kasten regelmäßig zur Überwachung von Strom- und Verbindungsstatus.
	\item Die Pumpe sollte geräuscharm arbeiten und keine Vibrationen erzeugen.
\end{itemize}

\textbf{Hinweis:} Sollte die blaue LED weiterhin blinken, ist keine Bluetooth-Verbindung aktiv. Überprüfen Sie in diesem Fall die App oder starten Sie den Verbindungsprozess erneut.